\documentclass[UTF8,zihao=-4,fontset=windows]{ctexart}
\usepackage[a4paper,margin=21mm,headheight=0pt,headsep=0pt,footskip=0pt]{geometry}
\usepackage{graphicx,booktabs,tabularx,array,longtable,xcolor,pdflscape}
\usepackage[unicode,colorlinks=true,linkcolor=black,urlcolor=blue]{hyperref}
\setCJKmainfont{Microsoft YaHei}
\setCJKsansfont{Microsoft YaHei}
\setmainfont{Times New Roman}
\definecolor{brand}{HTML}{4F46E5}
\ctexset{section={format=\Large\bfseries\color{brand}},subsection={format=\large\bfseries}}
\pagestyle{empty}
\setlength{\parindent}{2em}
\setlength{\parskip}{4pt}
\renewcommand{\arraystretch}{1.28}
\setlength{\tabcolsep}{5pt}
\hypersetup{pdftitle={实验二 软件设计与建模：活动票务与报名系统},pdfauthor={万连勇 241820015},pdfsubject={完整AI设计稿，待本人最终审查}}
\begin{document}
\begin{titlepage}
\centering\vspace*{25mm}
{\small SOFTWARE ENGINEERING · LAB 02\par}
\vspace{12mm}
{\Huge\bfseries 实验二\par}\vspace{5mm}
{\LARGE 软件设计与建模\par}\vspace{13mm}
{\Large\color{brand} 活动票务与报名系统\par}
\vspace{22mm}
\begin{tabular}{rl}
姓名：&万连勇\\[3mm]
学号：&241820015\\[3mm]
需求访谈对象：&桂鑫铭\\[3mm]
设计整理日期：&2026年11月13日\\[3mm]
版本：&完整AI设计稿，待本人最终审查
\end{tabular}
\vfill
\begin{minipage}{0.86\linewidth}\small
本文依据实验一最终需求进行设计，提供11个逻辑模块、六种UML、29张界面图及可编辑源文件。当前完成设计级核对与交互原型检查，尚未完成真实后端实现及其验收；人工思考、修改和Git记录由本人审查后据实补充。
\end{minipage}
\end{titlepage}
\setcounter{tocdepth}{1}
\pagenumbering{Roman}\tableofcontents\clearpage\pagenumbering{arabic}

\section{项目概况与主要功能}

\subsection{问题、目标与范围}
本项目面向活动组织与票务报名，解决活动发布、名额争抢、实名重复报名、候补递补及入场核验之间规则分散、状态不一致的问题。系统同时支持免费报名与收费票务，目标是使参与者获得清晰的报名结果，使组织者可维护活动和统计名单，使检票人员按范围核验，使管理员完成审核、授权、异常处理和追溯。

范围包括活动—场次—票种—分区座位、多人整单报名、模拟支付、15分钟占位、候补、退票、转赠、在线与离线检票、缺席纠错、取消延期、通知、导出统计和审计。当前实验只交付设计与界面原型，不实现生产软件。支付、退款、微信登录和订阅推送按实验一约定模拟；不增加真实资金渠道或身份认证服务。


\subsection{需求依据与版本}
需求基线为公开仓库software\_project的lab1-final标签，共40条需求：32条功能需求、6条非功能需求、2条约束。沿用已确认业务规则，基线中仍待细化的内容保持待澄清，不因AI提出设计建议而自动变为甲方确认。

\noindent 仓库：\url{https://github.com/wly-nju/software_project}\par
\noindent 需求基线提交：\texttt{ad83340e2cab6b91488b7a5732da7141c9088e41}\par
实验二工作开始提交为e7d1c0f。报告生成时完整设计稿尚未提交，当前不虚填实验二最终提交编号；提交后可在版本记录中补充。


\subsection{角色与权限}
\begin{small}
\noindent\begin{tabularx}{\linewidth}{p{24mm}X}
\toprule
角色 & 允许操作与边界 \\ \midrule
参与者 & 公开查询、实名报名、本人订单/候补/通知、合法取消退票和转赠；不能管理他人记录。 \\
组织者 & 仅管理授权活动，编辑、申请审核/取消/延期、名单导出和统计、误核销更正申请；不能直接批准取消或设置管理员例外比例。 \\
检票人员 & 仅核验授权场次，离线记录及恢复同步；不展示身份证和电话，不提供批量导出。 \\
平台管理员 & 活动审批、授权、取消及退款批次、纠错复核、审计查询；管理员身份不能兼任普通身份。 \\
\bottomrule
\end{tabularx}
\end{small}\par\smallskip

\clearpage

\section{模块划分与设计取舍}
采用11个逻辑模块表达职责。它们运行于同一后端应用，模块边界用于组织类、接口和规则，不等同于11个独立服务。权限、安全和审计贯穿所有受保护操作。

\begin{small}
\noindent\begin{tabularx}{\linewidth}{p{47mm}X}
\toprule
模块 & 主要职责 \\ \midrule
M01 用户与身份权限 & 账号登录、模拟微信登录、观演人资料、活动/场次授权、数据访问与字段可见性校验 \\
M02 活动场次与票种配置 & 活动、场次、票种和座位配置；整体审核、计划发布、下架；核心变更审核、容量调整、停售；维护取消和延期的活动状态 \\
M03 活动查询与开售提醒 & 多条件查询、排序、详情、状态及价格区间展示，开售提醒订阅；读取当前可售结果 \\
M04 订单票据与占位 & 主订单及子票单、下单快照、实名与座位唯一占位、幂等提交、模拟支付、15分钟超时；统一维护票据当前状态和凭证有效性 \\
M05 候补与自动递补 & 按观演人/场次/票种入队与退出、持久化顺位、释放名额优先递补、截止失效；收费递补接入 M04 支付期 \\
M06 退票退款与取消联动 & 免费取消、普通及延期退票、模拟退款、取消范围内订单/候补/票据联动；进行中取消的退款比例设定、批次冻结、失败重试 \\
M07 电子票转赠 & 场次开关、发起/接收/撤回/拒绝/过期、单次向外转赠与原路转回；协调票据锁定、持票人和凭证更换 \\
M08 检票缺席与纠错 & 授权场次核销、重复核销提示、离线记录与同步、冲突复核、结束后缺席认定、更正申请和审批 \\
M09 业务通知 & 站内消息、关联业务结果、模拟微信订阅推送；为开售、递补、变更、取消及退票提供通知 \\
M10 名单与容量统计 & 按授权范围多条件筛选及导出名单；按子票单统计容量、报名、候补、检票、缺席等数据 \\
M11 审计与安全记录 & 接收关键业务、安全及自动规则操作日志；按人/时间/类型/对象查询归档；禁止编辑删除、禁止记录完整敏感信息 \\
\bottomrule
\end{tabularx}
\end{small}\par\smallskip

\subsection{为何将票据状态集中维护}
M04负责订单、子票单、有效占位与当前凭证，M05、M06、M07和M08通过应用用例发起受控状态转换。退款模块不能绕过核销状态，转赠模块不能独立改持票人后遗漏凭证，检票端不能只相信旧二维码。这样的组织方式减少跨模块分别更新同一资格导致的冲突。

M02维护活动与容量配置，M06协调变更造成的票务后果。M09通知与M11审计分开：通知面向用户，审计面向追溯；关键审计失败时业务回滚，通知在业务提交后处理，不因通知迟到延长支付期限。M10统一汇总口径，避免列表、订单和报表分别定义库存。


\clearpage

\section{整体架构与数据设计}

\subsection{分层与部署选择}
采用模块化单体与单个MySQL数据库。界面层处理输入与反馈；应用层校验权限、组织用例和短事务；领域层通过实体与服务维护不变量；基础设施层提供持久化、模拟适配器、加密及离线缓存。课程原型共享占位和票据状态，本地事务比引入分布式服务更便于明确原子性和恢复边界。后端框架可在实现阶段选择。

所有影响同场次资格的事务先锁场次，再按稳定顺序锁票种、订单和票据。此设计用部分并发度换取清晰的一致性边界，是否满足目标QPS须实测；若锁等待成为瓶颈再依据结果缩小锁粒度。模拟支付和退款调用不持有数据库事务，结果返回后另开短事务应用。

组件图展示主要依赖与调用方向，应用用例入口协调权限、审计和提交。关键业务与审计、待通知事件同事务提交，M09在提交后发送站内消息和模拟推送。数据访问通过仓储与工作单元组织，界面不直接修改数据库。

\begin{landscape}
\subsection{组件图：11个模块与主要交互}
\begin{center}\includegraphics[width=\linewidth,height=0.75\textheight,keepaspectratio]{\detokenize{../UML/01_组件图.png}}\end{center}
\end{landscape}
\begin{landscape}
\subsection{类图：核心票务领域}
\begin{center}\includegraphics[width=\linewidth,height=0.75\textheight,keepaspectratio]{\detokenize{../UML/04_类图_核心票务.png}}\end{center}
Event组合Session，Session组合TicketType与SessionSeat；Order组合1～4张Ticket。Ticket引用当前实名身份、当前持票人及价格规则快照，以ActiveClaim表达当前占位，以Credential表达不可复用的历史凭证版本。TicketService组织创建、支付、超时及受控状态转换。

\end{landscape}
\begin{landscape}
\subsection{类图：身份授权与可追溯记录}
\begin{center}\includegraphics[width=\linewidth,height=0.69\textheight,keepaspectratio]{\detokenize{../UML/05_类图_身份与业务记录.png}}\end{center}
角色采用RoleGrant关联用户与业务范围，避免用继承固化角色兼任。支付关联原账户，退款关联子票及可选批次；转赠、核销、纠错、通知和审计分别留存，当前票据失效不覆盖曾经检票、曾经转赠等事实。

\end{landscape}

\subsection{存储关联、唯一约束与历史}
\begin{small}
\noindent\begin{tabularx}{\linewidth}{p{27mm}X}
\toprule
对象 & 关键约束与同步规则 \\ \midrule
账号与身份 & 账号唯一；证件和电话加密存储。规范化身份HMAC键用于同身份证跨账号唯一校验，密钥与数据库分开。 \\
主单与子票 & 请求号唯一并绑定调用者、载荷；成交价格与规则保存快照；同一主单只含一个场次，最多4张。 \\
有效占位 & 唯一(场次,身份键)、唯一(场次,非空座位)、每票至多一条。历史票可保留相同身份，不能用永久唯一历史索引阻止重报。 \\
凭证与转赠 & 凭证摘要唯一、每票最多一个当前有效凭证；旧凭证永久失效。待接收申请最多一条，向外成功一次，返还不重置次数。 \\
队列与流水 & 有效候补(身份,场次,票种)唯一，服务端序号持久化；支付流水、退款键和设备事件号分别唯一。 \\
当前与历史 & 状态、占位、持票人及新旧凭证同步更新；保留原订单、当前票据、退款流水、核销事实、更正与审计关联。 \\
\bottomrule
\end{tabularx}
\end{small}\par\smallskip
金额采用精确十进制与币种，不采用二进制浮点；具体舍入策略仍待确认。数据库事务或唯一约束失败时整单回滚，不能只依赖按钮防重复。敏感字段不能由前端遮挡代替数据库保护；日志不保存完整证件、电话、凭证或密钥。


\subsection{容量守恒与统计}
按子票单计数，总容量C＝有效报名E＋待支付L＋退票中R＋尚未分配U。E包含已报名、已检票和已缺席；候补和开售提醒不占位。示例100＝80＋10＋5＋5。U不是必然允许直接购买的库存，还受候补优先、停售及取消限制。容量不能调至E＋L＋R以下，已占座位不能移除。


\clearpage

\section{角色用例与状态约束}

\subsection{用例图：参与者与管理边界}
\noindent\begin{minipage}[t]{0.54\linewidth}\vspace{0pt}\begin{center}\includegraphics[width=\linewidth,height=0.55\textheight,keepaspectratio]{\detokenize{../UML/02_用例图_参与者.png}}\end{center}
\end{minipage}\hfill\begin{minipage}[t]{0.43\linewidth}\vspace{0pt}\begin{center}\includegraphics[width=\linewidth,height=0.55\textheight,keepaspectratio]{\detokenize{../UML/03_用例图_管理与检票.png}}\end{center}
\end{minipage}
参与者继承访客公开浏览能力；报名包含实名、库存和时段校验。管理员、组织者与检票人员的入口按权限分开，组织者提交申请，管理员审批和处理例外。超时、递补和缺席由系统自动规则执行，不虚构一个人工“系统用户”替代业务权限。


\clearpage

\subsection{状态图：子票单与受控例外}
\begin{center}\includegraphics[width=\linewidth,height=0.69\textheight,keepaspectratio]{\detokenize{../UML/06_状态图_子票单.png}}\end{center}
收费票创建待支付，免费票直接已报名。允许待支付→已报名/已取消，已报名→已检票/已缺席/退票中/合法取消，退票中→已退票或活动正常时失败恢复已报名；取消场次使相关资格失效。已退票、已取消和取消场次失效不得普通复活。经审批误操作更正与合法离线迟到是明确例外，不能绕过检票时段或取消状态。转赠待接收是正交锁，不新增票据生命周期状态。


\clearpage

\section{关键业务交互与冲突处理}

\subsection{时序图：整单创建、幂等与支付}
\begin{center}\includegraphics[width=\linewidth,height=0.75\textheight,keepaspectratio]{\detokenize{../UML/07_时序图_整单下单与支付.png}}\end{center}
同场次最多4张可跨票种，所有观演人和座位一起校验；任一项冲突整单失败。相同请求返回原订单，支付流水只生效一次。区分实际付款时间与通知到达；已关闭释放后收到支付结果只进入异常退款，不恢复票据。


\clearpage

\subsection{时序图：超时释放与候补优先}
\begin{center}\includegraphics[width=\linewidth,height=0.75\textheight,keepaspectratio]{\detokenize{../UML/08_时序图_超时与候补.png}}\end{center}
普通收费和收费递补均保留15分钟支付期，报名截止不缩短原期限。释放与同场同票种有效候补递补在协调事务内处理，防止普通下单插队；不跨票种，成功后清理同场其他候补。具体候补选座算法和恰好15分钟竞态仍待确认。


\clearpage

\subsection{时序图：转赠锁定、实名与凭证替换}
\begin{center}\includegraphics[width=\linewidth,height=0.75\textheight,keepaspectratio]{\detokenize{../UML/09_时序图_转赠与检票竞争.png}}\end{center}
转赠只允许已支付、有效、未检票且场次开启的票。期限取发起24小时与开演前30分钟较早者；截止后不能发起或接收。待接收锁定核销、退票及再次转赠；接收方实名且同场无有效占位。成功时持票人、身份占位、历史和凭证一起变更；向外一次，仅能转回原购票人，原订单和容量不变。


\clearpage

\subsection{时序图：取消标志与冻结退款批次}
\begin{center}\includegraphics[width=\linewidth,height=0.75\textheight,keepaspectratio]{\detokenize{../UML/10_时序图_取消与退款批次.png}}\end{center}
先提交取消标志，立即拒绝下单、递补、核销和转赠，再幂等处理关联记录及退款。进行中因主办方取消已检票付费默认全退，管理员执行前可附理由设0\%～100\%，启动后冻结；重试沿用原键、比例和金额。已结束后取消的已检票/已缺席不退款；所有退款回原支付账户，失败不恢复票据。


\clearpage

\section{部署与质量设计}

\subsection{部署图：本地单实例与离线端}
\begin{center}\includegraphics[width=\linewidth,height=0.47\textheight,keepaspectratio]{\detokenize{../UML/11_部署图.png}}\end{center}
在线端经本地HTTP调用单实例后端和MySQL，生产及小程序环境按需求使用HTTPS。检票端可有离线缓存和设备事件队列；断网无法实时获知撤权、取消、转赠或其他设备核销，界面明确提示风险。联网后按设备事件号去重、核对当前状态，冲突交管理员复核，取消票据不因上传恢复资格。离线资料下载、有效期和清除策略保持待细化。


\subsection{性能计划与判定}
\begin{small}
\noindent\begin{tabularx}{\linewidth}{Xrrr}
\toprule
接口 & QPS/每轮 & 平均ms & P99 ms \\ \midrule
列表 & 50 / 1000 & ≤500 & ≤1000 \\
详情 & 50 / 1000 & ≤800 & ≤1500 \\
下单 & 20 / 500 & ≤1500 & ≤3000 \\
核销 & 30 / 800 & ≤500 & ≤1000 \\
\bottomrule
\end{tabularx}
\end{small}\par\smallskip
沿用8核、16GB、SSD、Windows11/macOS14、本地单实例与MySQL；不少于100活动×2场次×3票种、10000报名、1000用户。预热100次，每组三轮；每轮P99取第ceil(0.99N)个最近秩样本，再求三轮指标均值。保留失败样本，全部正式请求成功率100\%，技术错误/超时/>3000ms/一致性错误计失败；正确业务拒绝及幂等返回算成功。需覆盖真实成功路径并核对数据库，不以快速拒绝冒充下单性能。


\clearpage

\subsection{安全、恢复、易用性与反馈}
\begin{small}
\noindent\begin{tabularx}{\linewidth}{p{31mm}X}
\toprule
需求 & 设计与后续验收 \\ \midrule
REQ-034/035 & 逐用例鉴权与归属过滤，约定越权100\%拒绝且业务不变；数据库密文、最少字段、导出与日志脱敏，生产HTTPS。 \\
REQ-036 & 请求、流水、设备事件及任务幂等，重启恢复未完成任务；按三类故障注入验证30秒恢复和数据一致性，不把99.9\%年度目标说成已实测。 \\
REQ-037 & 购票核心4步；首次用户至少4/5在3分钟内完成，组织者至少3/5在8分钟内完成。未完成者平均耗时口径仍待确认。 \\
REQ-038 & 核心操作3秒内显示服务端业务反馈；本地加载、处理中与成功分别表达，不能凭原型提示声称接口已经受理。 \\
REQ-032/039 & 关键日志与业务同事务；状态转换受控、历史保留、日志不可编辑删除；权限撤销、取消和转赠竞争需后端验证。 \\
\bottomrule
\end{tabularx}
\end{small}\par\smallskip

\subsection{重点验收场景与当前限制}
后续实现应覆盖最后名额竞争、多人整单中一张冲突、同请求/流水重试、支付超时竞争、候补优先、普通退票失败恢复、取消时退款失败、转赠与核销互斥、离线重复与迟到、纠错审批、审计故障回滚以及容量守恒。当前只有设计级检查和原型交互检查，没有真实后端测试、性能结果或人测数据。

沿用基线待澄清项：密码和账号找回、输入上限、过期发布计划、混合取消状态聚合、重复提醒计数、模拟失败/取消支付后的政策、15分钟精确端点、金额舍入、候补选座、检票结束端点、离线缓存策略、延期退票截止/已检票适用范围、通知重试与接收对象、导出格式/上限及易用性未完成样本口径。实现前应确认，不由AI写成既定业务。


\clearpage

\section{UI界面与功能说明}

\subsection{工具、信息架构与视觉规则}
UI以diagrams.net可编辑mxGraph格式组织29页画板，文字、容器和控件可分别编辑；HTML/CSS/JavaScript提供交互补充，本机Chrome导出PNG。源文件与原型均保留，未使用或冒称Figma/Pixso。PNG为主预览，draw.io交换稿打开后的字体与细节可能需要微调。

参与者使用移动布局，底部导航连接活动、票夹、候补和我的；组织者与管理员使用侧栏工作台，检票端采用大反馈区。紫色表示主操作、绿色表示有效、黄色表示待处理、红色表示风险，同时提供文字状态。菜单展示各角色页面便于审阅，真实权限须在入口和服务端落实。

购票路径为票种数量→观演人→确认信息→提交支付；注册、首次登录、浏览和支付渠道确认不混入核心4步口径。转赠锁定、离线核验风险、售罄、处理中、403和取消退款失败均有独立反馈。资料为测试数据，二维码是不可核销的装饰示例。多人混合票种控件可演示计价，支付导航使用默认单张测试订单。

\begin{small}
\noindent\begin{tabularx}{\linewidth}{lX}
\toprule
模块 & 界面编号 \\ \midrule
M01 & 01、04、12、23、26、27 \\
M02 & 03、13、14、15、16、17、25、28 \\
M03 & 02、03、26 \\
M04 & 04、05、06、07、15、25、26、29 \\
M05 & 08 \\
M06 & 06、09、17、28 \\
M07 & 07、10、15、26、29 \\
M08 & 18、19、20、26 \\
M09 & 11 \\
M10 & 13、21、22 \\
M11 & 24 \\
\bottomrule
\end{tabularx}
\end{small}\par\smallskip
下列画面逐页说明全部29张UI。移动页每组2张，管理页每页1张并采用横向版面，优先保证文字可读；独立PNG可放大查看，完整需求映射见附录与设计说明。原型行为不等于后端真实执行。


\clearpage

\subsection{移动界面：01 / 02}
\noindent\begin{minipage}[t]{0.48\linewidth}\vspace{0pt}\centering\begin{center}\includegraphics[width=\linewidth,height=0.72\textheight,keepaspectratio]{\detokenize{../UI界面/01_login_登录与身份入口.png}}\end{center}
{\small\bfseries 01 登录与身份入口}\par\end{minipage}
\hfill
\begin{minipage}[t]{0.48\linewidth}\vspace{0pt}\centering\begin{center}\includegraphics[width=\linewidth,height=0.72\textheight,keepaspectratio]{\detokenize{../UI界面/02_explore_活动查询.png}}\end{center}
{\small\bfseries 02 活动查询}\par\end{minipage}

\vspace{4mm}
\noindent\textbf{01 M01}：账号登录及模拟微信授权；普通用户入口不提供自行授予管理员权限。\par
\noindent\textbf{02 M03}：关键词、城市、分类、日期、价格及状态组合筛选；公开可见不等于允许购买。\par

\clearpage

\subsection{移动界面：03 / 04}
\noindent\begin{minipage}[t]{0.48\linewidth}\vspace{0pt}\centering\begin{center}\includegraphics[width=\linewidth,height=0.72\textheight,keepaspectratio]{\detokenize{../UI界面/03_detail_活动详情与提醒.png}}\end{center}
{\small\bfseries 03 活动详情与提醒}\par\end{minipage}
\hfill
\begin{minipage}[t]{0.48\linewidth}\vspace{0pt}\centering\begin{center}\includegraphics[width=\linewidth,height=0.72\textheight,keepaspectratio]{\detokenize{../UI界面/04_checkout_实名整单确认.png}}\end{center}
{\small\bfseries 04 实名整单确认}\par\end{minipage}

\vspace{4mm}
\noindent\textbf{03 M02/M03}：场次、价格、售票状态和退票规则；开售提醒与候补分别设置。\par
\noindent\textbf{04 M01/M04}：选择最多4个观演人，整单校验并占位，展示成交价及规则快照。\par

\clearpage

\subsection{移动界面：05 / 06}
\noindent\begin{minipage}[t]{0.48\linewidth}\vspace{0pt}\centering\begin{center}\includegraphics[width=\linewidth,height=0.72\textheight,keepaspectratio]{\detokenize{../UI界面/05_payment_模拟支付与超时.png}}\end{center}
{\small\bfseries 05 模拟支付与超时}\par\end{minipage}
\hfill
\begin{minipage}[t]{0.48\linewidth}\vspace{0pt}\centering\begin{center}\includegraphics[width=\linewidth,height=0.72\textheight,keepaspectratio]{\detokenize{../UI界面/06_orders_订单列表与状态.png}}\end{center}
{\small\bfseries 06 订单列表与状态}\par\end{minipage}

\vspace{4mm}
\noindent\textbf{05 M04}：15分钟占位；可演示成功、失败及原订单恢复，不模拟真实资金收付。\par
\noindent\textbf{06 M04/M06}：原订单恢复、免费取消和单张退票入口，历史记录不复活。\par

\clearpage

\subsection{移动界面：07 / 08}
\noindent\begin{minipage}[t]{0.48\linewidth}\vspace{0pt}\centering\begin{center}\includegraphics[width=\linewidth,height=0.72\textheight,keepaspectratio]{\detokenize{../UI界面/07_ticket_电子票与锁定状态.png}}\end{center}
{\small\bfseries 07 电子票与锁定状态}\par\end{minipage}
\hfill
\begin{minipage}[t]{0.48\linewidth}\vspace{0pt}\centering\begin{center}\includegraphics[width=\linewidth,height=0.72\textheight,keepaspectratio]{\detokenize{../UI界面/08_wait_候补登记与递补.png}}\end{center}
{\small\bfseries 08 候补登记与递补}\par\end{minipage}

\vspace{4mm}
\noindent\textbf{07 M04/M07}：当前持票、原订单、凭证状态及转赠锁定；二维码仅为不可核销的视觉示例。\par
\noindent\textbf{08 M05}：独立候补序号及跨票种排队；成功后清理同场其他候补；付费递补原15分钟期限。\par

\clearpage

\subsection{移动界面：09 / 10}
\noindent\begin{minipage}[t]{0.48\linewidth}\vspace{0pt}\centering\begin{center}\includegraphics[width=\linewidth,height=0.72\textheight,keepaspectratio]{\detokenize{../UI界面/09_refund_退票申请与进度.png}}\end{center}
{\small\bfseries 09 退票申请与进度}\par\end{minipage}
\hfill
\begin{minipage}[t]{0.48\linewidth}\vspace{0pt}\centering\begin{center}\includegraphics[width=\linewidth,height=0.72\textheight,keepaspectratio]{\detokenize{../UI界面/10_transfer_转赠发起与接收.png}}\end{center}
{\small\bfseries 10 转赠发起与接收}\par\end{minipage}

\vspace{4mm}
\noindent\textbf{09 M06}：按子票申请，显示手续费、原账户及状态；转赠待接收禁止申请。\par
\noindent\textbf{10 M07}：展示较早期限、单次向外转赠、接收资格及原路返还，不允许转第三人。\par

\clearpage

\subsection{移动界面：11 / 12}
\noindent\begin{minipage}[t]{0.48\linewidth}\vspace{0pt}\centering\begin{center}\includegraphics[width=\linewidth,height=0.72\textheight,keepaspectratio]{\detokenize{../UI界面/11_messages_消息中心.png}}\end{center}
{\small\bfseries 11 消息中心}\par\end{minipage}
\hfill
\begin{minipage}[t]{0.48\linewidth}\vspace{0pt}\centering\begin{center}\includegraphics[width=\linewidth,height=0.72\textheight,keepaspectratio]{\detokenize{../UI界面/12_spectators_观演人与实名资料.png}}\end{center}
{\small\bfseries 12 观演人与实名资料}\par\end{minipage}

\vspace{4mm}
\noindent\textbf{11 M09}：业务关联通知、变化原因与原期限；本人与他人数据隔离。\par
\noindent\textbf{12 M01}：代他人购票与个人资料最少展示；示例身份不代表真实实名认证渠道。\par

\clearpage

\subsection{移动界面：18 / 19}
\noindent\begin{minipage}[t]{0.48\linewidth}\vspace{0pt}\centering\begin{center}\includegraphics[width=\linewidth,height=0.72\textheight,keepaspectratio]{\detokenize{../UI界面/18_checkin_检票核销与重复提示.png}}\end{center}
{\small\bfseries 18 检票核销与重复提示}\par\end{minipage}
\hfill
\begin{minipage}[t]{0.48\linewidth}\vspace{0pt}\centering\begin{center}\includegraphics[width=\linewidth,height=0.72\textheight,keepaspectratio]{\detokenize{../UI界面/19_offline_离线核验与同步复核.png}}\end{center}
{\small\bfseries 19 离线核验与同步复核}\par\end{minipage}

\vspace{4mm}
\noindent\textbf{18 M08}：当前凭证、授权场次与检票时间；重复核销展示首次记录。\par
\noindent\textbf{19 M08}：显式离线风险及待上传记录；上传去重和异常交管理员处理。\par

\clearpage

\subsection{移动界面：25 / 27}
\noindent\begin{minipage}[t]{0.48\linewidth}\vspace{0pt}\centering\begin{center}\includegraphics[width=\linewidth,height=0.72\textheight,keepaspectratio]{\detokenize{../UI界面/25_seats_参与者分区选座.png}}\end{center}
{\small\bfseries 25 参与者分区选座}\par\end{minipage}
\hfill
\begin{minipage}[t]{0.48\linewidth}\vspace{0pt}\centering\begin{center}\includegraphics[width=\linewidth,height=0.72\textheight,keepaspectratio]{\detokenize{../UI界面/27_register_账号注册.png}}\end{center}
{\small\bfseries 27 账号注册}\par\end{minipage}

\vspace{4mm}
\noindent\textbf{25 M02/M04}：当前可选与已占座位区分，最终仍由服务端处理最后座位竞争。\par
\noindent\textbf{27 M01}：建立普通账号；管理员不通过公开注册取得身份，密码策略待确认。\par

\clearpage

\subsection{移动界面：29}
\noindent\begin{minipage}[t]{0.48\linewidth}\vspace{0pt}\centering\begin{center}\includegraphics[width=\linewidth,height=0.72\textheight,keepaspectratio]{\detokenize{../UI界面/29_转赠待接收状态.png}}\end{center}
{\small\bfseries 29 转赠待接收状态}\par\end{minipage}

\vspace{4mm}
\noindent\textbf{29 M04/M07}：票据锁定时明确禁止退票、核销和再次转赠；撤回后恢复有效状态。\par

\clearpage
\begin{landscape}
\subsection{管理界面：13 组织者活动工作台}
\begin{center}\includegraphics[width=0.98\linewidth,height=0.77\textheight,keepaspectratio]{\detokenize{../UI界面/13_organizer_组织者活动工作台.png}}\end{center}
\noindent\textbf{M02/M10 / 组织者}：活动状态、核心变更暂停和授权范围，提供完整管理入口。\par\end{landscape}

\clearpage
\begin{landscape}
\subsection{管理界面：14 活动创建与核心修改}
\begin{center}\includegraphics[width=0.98\linewidth,height=0.77\textheight,keepaspectratio]{\detokenize{../UI界面/14_editor_活动创建与核心修改.png}}\end{center}
\noindent\textbf{M02 / 组织者}：分步维护、草稿保存、整体提交审核；非核心文案与核心规则区分。\par\end{landscape}

\clearpage
\begin{landscape}
\subsection{管理界面：15 场次票种容量及座位}
\begin{center}\includegraphics[width=0.98\linewidth,height=0.77\textheight,keepaspectratio]{\detokenize{../UI界面/15_session_场次票种容量及座位.png}}\end{center}
\noindent\textbf{M02/M04/M07 / 组织者}：容量下限及区域合计校验、停售保留历史、按场次转赠开关。\par\end{landscape}

\clearpage
\begin{landscape}
\subsection{管理界面：16 管理员审核与发布控制}
\begin{center}\includegraphics[width=0.98\linewidth,height=0.77\textheight,keepaspectratio]{\detokenize{../UI界面/16_review_管理员审核与发布控制.png}}\end{center}
\noindent\textbf{M02 / 管理员}：活动整体审核、驳回原因、下架；组织者无审批权限。\par\end{landscape}

\clearpage
\begin{landscape}
\subsection{管理界面：17 取消审批与退款批次}
\begin{center}\includegraphics[width=0.98\linewidth,height=0.77\textheight,keepaspectratio]{\detokenize{../UI界面/17_cancel_取消审批与退款批次.png}}\end{center}
\noindent\textbf{M06/M02 / 管理员}：默认全额、执行前比例调整、执行后冻结、原路退款、失败重试与历史保留。\par\end{landscape}

\clearpage
\begin{landscape}
\subsection{管理界面：20 核销异常与纠错审批}
\begin{center}\includegraphics[width=0.98\linewidth,height=0.77\textheight,keepaspectratio]{\detokenize{../UI界面/20_correction_核销异常与纠错审批.png}}\end{center}
\noindent\textbf{M08 / 管理员}：组织者附原因申请，管理员批准；普通终态不能借纠错复活。\par\end{landscape}

\clearpage
\begin{landscape}
\subsection{管理界面：21 名单筛选与脱敏导出}
\begin{center}\includegraphics[width=0.98\linewidth,height=0.77\textheight,keepaspectratio]{\detokenize{../UI界面/21_roster_名单筛选与脱敏导出.png}}\end{center}
\noindent\textbf{M10 / 组织者}：场次、票种、状态及时间组合筛选，授权范围导出并审计，敏感字段脱敏。\par\end{landscape}

\clearpage
\begin{landscape}
\subsection{管理界面：22 容量与参与情况统计}
\begin{center}\includegraphics[width=0.98\linewidth,height=0.77\textheight,keepaspectratio]{\detokenize{../UI界面/22_statistics_容量与参与情况统计.png}}\end{center}
\noindent\textbf{M10 / 组织者}：容量守恒、按子票计数、预约提醒与候补分开、历史统计与当前可售区分。\par\end{landscape}

\clearpage
\begin{landscape}
\subsection{管理界面：23 角色与范围授权}
\begin{center}\includegraphics[width=0.98\linewidth,height=0.77\textheight,keepaspectratio]{\detokenize{../UI界面/23_roles_角色与范围授权.png}}\end{center}
\noindent\textbf{M01 / 管理员}：普通角色兼任、管理员隔离、按活动或场次授权，旧页面不能绕过当前服务端权限。\par\end{landscape}

\clearpage
\begin{landscape}
\subsection{管理界面：24 审计日志查询}
\begin{center}\includegraphics[width=0.98\linewidth,height=0.77\textheight,keepaspectratio]{\detokenize{../UI界面/24_audit_审计日志查询.png}}\end{center}
\noindent\textbf{M11 / 管理员}：按人、时间、动作、对象查询，展示前后状态和原因，不提供修改删除。\par\end{landscape}

\clearpage
\begin{landscape}
\subsection{管理界面：26 异常与边界反馈}
\begin{center}\includegraphics[width=0.98\linewidth,height=0.77\textheight,keepaspectratio]{\detokenize{../UI界面/26_states_异常与边界反馈.png}}\end{center}
\noindent\textbf{M01/M03/M04/M07/M08 / 管理员}：统一展示空结果、售罄、重复请求、权限拒绝、锁定及取消退款失败等边界状态。\par\end{landscape}

\clearpage
\begin{landscape}
\subsection{管理界面：28 组织者取消延期申请}
\begin{center}\includegraphics[width=0.98\linewidth,height=0.77\textheight,keepaspectratio]{\detokenize{../UI界面/28_change_request_组织者取消延期申请.png}}\end{center}
\noindent\textbf{M02/M06 / 组织者}：组织者提交取消或延期申请，不能直接批准或设定管理员退款例外比例。\par\end{landscape}

\clearpage

\section{AI辅助使用过程与修改优化}
本章按实际辅助工作记录。工具为Codex桌面端，本轮未获得可核验的具体模型版本，因此不沿用实验一型号。输入为实验一lab1-final需求、最新版实验二任务书、现有模块初稿及本轮“先完成所有内容，再由本人审查”的指令。AI产生的思路、检查及修改在此明确标识，尚未由本人作出最终取舍。

\begin{small}
\noindent\begin{tabularx}{\linewidth}{p{17mm}p{22mm}X}
\toprule
记录 & 场景 & 结果、优缺点与优化 \\ \midrule
AI-001 & 需求映射 & 11模块与40条追踪初稿。优点：覆盖清晰；局限：尚无事务/界面证据。保留11月11日原始快照，本轮继续细化。 \\
AI-002 & 架构数据 & 模块化单体、类与关联、唯一约束、场次锁、短事务、取消标志与批次。优点：一致性边界可核对；局限：锁性能与密钥/缓存策略待验证。 \\
AI-003 & PlantUML & 生成11张六种图并实际渲染。初始时序换行过密，调大消息宽度；业务类图缩短注释换行。保留首次图形结果用于比较。 \\
AI-004 & UI设计 & 移动/管理/检票布局与29张画面。补充独立注册、组织者取消延期申请，避免普通用户进入管理员操作；补多人混合票种控件、转赠锁定状态和无浮层截图。 \\
AI-005 & 报告整理 & 汇总架构理由、数据状态、全部UML/UI和本章，保留LaTeX。优点：便于图文追踪；局限：版式和个人体验仍需核对，不能制造人工评价。 \\
AI-006 & 设计核查 & 源文件配对、40条追踪、11模块覆盖、原型7项检查、PDF渲染与打包。只形成设计级证据，不宣称并发/安全/性能验收通过。 \\
\bottomrule
\end{tabularx}
\end{small}\par\smallskip

\clearpage

\subsection{AI自检中实际发现并修改的内容}
\begin{small}
\noindent\begin{tabularx}{\linewidth}{p{45mm}X}
\toprule
原稿问题 & 修改与原因 \\ \midrule
时序图消息过度折行 & 将PlantUML消息宽度由50调整为260，降低图高，增加报告可读性；不改变业务内容。 \\
首次登录缺独立注册页 & 增加普通账号注册，禁止公开自行取得管理员身份，密码策略待确认。 \\
组织者工作台指向管理员取消处理 & 增加组织者取消/延期申请页，工作台指向申请入口；执行取消及例外比例仍在管理员页。 \\
单人确认示例不足以展示多人混票 & 补4人选择和各自票种、动态计价控件；明确支付演示用默认1张，不把原型当完整购物车。 \\
转赠状态截图残留提示浮层 & 保留稳定的锁定票据画面，便于检查禁止退票/核销的状态文案。 \\
\bottomrule
\end{tabularx}
\end{small}\par\smallskip
上述修改为AI对自己输出的核对和优化，不计作本人已完成的人工审查。本人应结合原始需求、课程范围和后续实现能力，判断是否接受集中状态维护、场次锁、离线风险提示与管理流程，记录真实的保留或修改理由，并用Git保留前后差异。原始快照、设计源文件及检查结果在工作目录中保留。


\subsection{辅助工具的作用与局限}
AI适合将长需求映射成模块、初步列出实体和冲突、生成PlantUML及重复界面骨架，并帮助检查编号和产物配对。局限是图能渲染不代表业务正确，界面可点击不代表接口存在，合理的技术选择也不能替代性能实测。尤其容易把待确认参数补成常识，或把组织者和管理员权限混用，因此需要保留来源和状态。


\subsection{可否让人只做检查校对}
可建立“冻结需求及未决项→AI生成带编号设计→自动检查覆盖、语法和界面→人检查业务取舍、权限、异常与成本→AI按意见修改→再次检查→Git保存可比版本”的流程。人能减少重复绘图和排版，但不能只做文字校对：目标冲突、风险接受、未决规则和实现成本仍须由人决定。自动检查能帮助定位问题，不能代替客户确认，也不能生成并不存在的人测和人工修改记录。


\clearpage

\section{实验总结与审查交接}
本轮设计的主要难点是票据状态、转赠锁、退款进度和核销历史的关系，以及下单、候补、取消对同一占位的竞争。设计通过独立有效占位、统一票据状态入口、事务与唯一约束、取消标志先行、冻结退款批次以及保留不可覆盖历史加以表达。界面同时暴露锁定、失败和离线风险，避免只呈现顺利流程。

已形成需求追踪、模块与数据设计、11张PlantUML及PNG、29张UI及可编辑画板/交互源文件、质量验收计划和本报告。设计级验证包括源文件实际渲染、页面脚本和横向布局检查、7项原型交互；它们仅说明当前交付可审阅，不能证明实际业务系统满足并发、安全、恢复和性能要求。


\subsection{本人最终审查时应完成的事项}
\begin{enumerate}
\item 阅读40条需求追踪，核对权限、整单、候补、取消退款、转赠与离线检票等关键规则。
\item 决定是否接受11模块边界、模块化单体、统一票据入口、身份键和场次锁设计；未决项继续标注。
\item 在人工审查记录中写实际判断、修改内容和依据，在本章补充自己的使用体验及取舍；正确内容可以保留，不强行制造错误。
\item 先保存完整AI设计稿Git提交，再修改相同源文件并保存人工修订提交；不移动实验一固定标签。
\item 修改后重新导图、导出受影响UI、编译并复核报告，重建提交压缩包，检查最终打包结构。
\end{enumerate}

\subsection{资料来源}
\begin{itemize}
\item 本仓库lab1-final:requirements.md及访谈确认记录。
\item 本仓库lab2/实验2要求.md（用户提供的最新版任务书）。
\item PlantUML类图官方文档：\url{https://plantuml.com/zh/class-diagram}
\item PlantUML组件图官方文档：\url{https://plantuml.com/zh/component-diagram}
\item PlantUML时序图官方文档：\url{https://plantuml.com/zh/sequence-diagram}
\end{itemize}

\clearpage

\section{附录：逐条需求与设计追踪}
UML编号对应文件名前两位，UI编号对应界面前两位。无独立界面的横切要求见架构和质量设计；本表只表明设计落点。

\begin{small}\begin{longtable}{p{17mm}p{40mm}p{16mm}p{20mm}p{45mm}}\toprule 需求 & 名称 & 主责 & UML & UI/落实方式\\\midrule\endhead
REQ-001 & 项目范围与外部渠道模拟 & 横切 & 01、11 & 01、05 \\
REQ-002 & 注册登录与身份授权 & M01 & 02、03、05 & 01、12、23、27 \\
REQ-003 & 按角色与归属限制业务访问 & M01 & 02、03、05 & 01、12、13、18、23、28 \\
REQ-004 & 活动场次与票种资料维护 & M02 & 01、04 & 03、13、14、15 \\
REQ-005 & 活动审核发布与下架 & M02 & 03、04 & 13、14、16 \\
REQ-006 & 已发布活动变更管理 & M02 & 01、03 & 13、14、16 \\
REQ-007 & 容量调整与票种停用 & M02 & 04、08 & 15、22 \\
REQ-008 & 多条件活动查询与列表展示 & M03 & 02、01 & 02、03、26 \\
REQ-009 & 开售提醒订阅 & M03 & 02、05 & 02、03、11 \\
REQ-010 & 多人多票种整单报名 & M04 & 04、07 & 04、12、25 \\
REQ-011 & 实名唯一占位与座位竞争 & M04 & 04、07 & 04、10、25、26 \\
REQ-012 & 重复提交与原订单恢复 & M04 & 07 & 04、05、06、26 \\
REQ-013 & 模拟支付及延迟结果处理 & M04 & 07 & 05 \\
REQ-014 & 待支付超时释放 & M04 & 08 & 05、08 \\
REQ-015 & 报名时段及当前状态校验 & M04 & 07 & 03、25 \\
REQ-016 & 免费报名取消与重报 & M06 & 06、08 & 06、09 \\
REQ-017 & 收费退票与模拟退款 & M06 & 06、08 & 06、09 \\
REQ-018 & 退票时限与金额计算 & M06 & 06 & 03、09 \\
REQ-019 & 候补申请与队列顺位 & M05 & 05、08 & 08 \\
REQ-020 & 释放名额自动递补 & M05 & 08 & 08 \\
REQ-021 & 在线二维码检票 & M08 & 03、09 & 07、18 \\
REQ-022 & 离线检票与恢复同步 & M08 & 03、11 & 19、20、26 \\
REQ-023 & 缺席认定与统计 & M08 & 06 & 19、20 \\
REQ-024 & 人工误操作更正审批 & M08 & 03、05、06 & 20 \\
REQ-025 & 活动或单场次取消联动 & M06 & 06、10 & 17、28 \\
REQ-026 & 活动延期持票与退票选择 & M06 & 01、03 & 09、17、28 \\
REQ-027 & 电子票转赠 & M07 & 05、09 & 07、09、10、15、26、29 \\
REQ-028 & 关键业务站内通知 & M09 & 01、05 & 11 \\
REQ-029 & 报名名单筛选与导出 & M10 & 03 & 21 \\
REQ-030 & 容量与参与情况统计 & M10 & 01、04 & 15、22 \\
REQ-031 & 业务关联与历史快照保存 & 横切 & 04、05 & 07 \\
REQ-032 & 关键操作日志与管理员查询 & M11 & 01、05 & 16、24 \\
REQ-033 & 指定负载下的接口性能 & 横切 & 07、11 & 02 \\
REQ-034 & 未授权访问防护 & 横切 & 02、03、11 & 23、24、26、27 \\
REQ-035 & 个人信息存储与展示保护 & 横切 & 05、11 & 12、21、24 \\
REQ-036 & 故障恢复与核心数据可靠性 & 横切 & 08、10、11 & 19 \\
REQ-037 & 首次使用者任务易用性 & 横切 & 02、03 & 04、14 \\
REQ-038 & 核心操作及时反馈 & 横切 & 07、09 & 26 \\
REQ-039 & 业务状态及允许转换约束 & 横切 & 06 & 06、07、18、29 \\
REQ-040 & 进行中取消的已检票人员退款政策 & M06 & 05、10 & 17 \\\bottomrule\end{longtable}\end{small}
\end{document}
